\subsection{Actual Gauss inflation and absence of isolated points on the whole line and integral curves}
\label{subsec:gauss-inflation-curve-nonisolation}

\begin{lemma}[A nonconstant actual connected path excludes an open singleton]
\label{topology:actual-connected-path-nonisolation}
\leanok
\lean{EnvelopeProperty.not_isOpen_singleton_of_continuous_connected}
Let $f:Z\to W$ be an actual continuous map from a preconnected
space to a $T_1$ space. If $z\in Z$ and some actual $z'\in Z$
has $f(z')\ne f(z)$, then $f(z)$ is nonisolated. No injection,
open map or separation stronger than $T_1$ is assumed.
\end{lemma}
\begin{proof}
If the image singleton were open, its actual inverse image would
be open. It is closed by $T_1$ and actual continuity, contains $z$,
and misses $z'$. This contradicts preconnectedness: every nonempty
clopen subset is the whole domain.
\end{proof}

\begin{definition}[An actual whole Gauss inflation of an original line point]
\label{spectrum:actual-line-gauss-inflation}
\leanok
\lean{EnvelopeProperty.Berkovich.RelativeSpectrum.lineGaussInflation,
EnvelopeProperty.Berkovich.RelativeSpectrum.lineGaussInflation_val_apply,
EnvelopeProperty.Berkovich.RelativeSpectrum.continuous_lineGaussInflation,
EnvelopeProperty.Berkovich.RelativeSpectrum.lineGaussInflation_zero,
EnvelopeProperty.Berkovich.RelativeSpectrum.radius_le_lineGaussInflation_val_X}
\uses{spectrum:centered-gauss-whole-ray}
Let $R$ be ANY commutative ring, $\nu$ ANY prescribed real-valued
valuation, and $x$ ANY actual point of the prescribed line spectrum
of $R[T]$. Regard $x$ as the genuine coefficient valuation on the
original ring $R[T]$. Form its actual translated Gauss seminorm on
$R[T][S]$ centered at the ORIGINAL element $T$, then pull it back
along the actual $R$-algebra map $R[T]\to R[T][S]$, $T\mapsto S$.
This constructs a whole prescribed line point $I_x(r)$ for every
$r\ge0$, continuously, with $I_x(0)=x$ and $r\le |T|_{I_x(r)}$.
No residue-field identification, field, completeness, norm bound,
support condition, point classification or path certificate is assumed.
\end{definition}
\begin{proof}
Both ring maps are the existing genuine polynomial algebra maps;
actual prescribed Gauss and pullback preserve EVERY original
scalar. Joint Gauss continuity followed by true fixed pullback
gives continuity in radius. At zero radius the centered Gauss
seminorm is whole evaluation $S\mapsto T$. Its composite with
$T\mapsto S$ is the IDENTITY on EVERY original polynomial,
proved by polynomial algebra-map extensionality. Thus the point
at radius zero is the ENTIRE original $x$.

The actual translation sends the outer variable to $S+T$.
Its coefficient of $S$ is the true unit. The genuine coefficient
lower bound for Gauss seminorms therefore gives
$r=|1|_x r\le |S+T|_{\mathrm{Gauss}(x,r)}=|T|_{I_x(r)}$.
This uses no informal tensor or residue-point comparison.
\end{proof}

\begin{theorem}[The whole actual line has no isolated points]
\label{spectrum:actual-all-line-points-nonisolation}
\leanok
\lean{EnvelopeProperty.Berkovich.RelativeSpectrum.not_isOpen_singleton_line,
EnvelopeProperty.Berkovich.RelativeSpectrum.neBot_nhdsNE_line}
\uses{topology:actual-connected-path-nonisolation,spectrum:actual-line-gauss-inflation,spectrum:topology,spectrum:relative}
EVERY actual point of the prescribed line spectrum over ANY
commutative ring and ANY real valuation is nonisolated. No
nonzero-support, rationality or residue-degree hypothesis is required.
\end{theorem}
\begin{proof}
The actual inflation path begins at $x$. At the genuine radius
$r=|T|_x+1$, its proved coordinate lower bound implies
$|T|_{I_x(r)}> |T|_x$, so this whole point differs from $x$.
The radius domain $\mathbb R_{\ge0}$ is genuinely connected and
the actual line spectrum is Hausdorff by the independently defined
evaluation topology. Apply the actual connected-path lemma.
The punctured-neighborhood identity yields filter nontriviality.
This is an actual nonisolation argument for the ENTIRE point,
without a support-density inference or classification of point types.
\end{proof}

\begin{theorem}[Every actual point of an integral affine curve is nonisolated]
\label{scheme:actual-all-integral-affine-curve-points-nonisolation}
\leanok
\lean{EnvelopeProperty.Berkovich.RelativeSchemeValuedPoint.not_isOpen_singleton_affine_of_dimension_le_one,
EnvelopeProperty.Berkovich.RelativeSchemeValuedPoint.neBot_nhdsNE_affine_of_dimension_le_one,
EnvelopeProperty.Berkovich.RelativeSchemeValuedPoint.not_isOpen_singleton_affine_of_dedekind}
\uses{algebra:actual-finite-free-integral-curve-line-parameter,spectrum:actual-all-line-points-nonisolation,spectrum:actual-finite-free-open-map,scheme:relative-affine-valued-point-homeomorphism}
Over ANY valued field, EVERY actual native point of a nonfield
finite-type integral affine curve of genuine dimension at most
one is nonisolated. In particular this applies to every actual
nonfield finite-type Dedekind domain over any field. No closed
support, perfectness, completeness or etale-chart hypothesis remains.
\end{theorem}
\begin{proof}
Use the ACTUAL finite free polynomial parameter constructed from
Noether normalization and genuine original prime heights. Its
whole seminorm restriction is genuinely open. Every image point
on the actual line is nonisolated by the preceding theorem; an
open source singleton would have an open image singleton.
The independently defined native affine homeomorphism retains
the ENTIRE original point and transfers this contradiction.
For Dedekind domains, the genuine dimension bound is derived
from the actual nonzero-prime maximality property. No analytic
potential or maximum-principle comparison is claimed.
\end{proof}
